% !TeX program = xelatex
\documentclass[11pt,a4paper]{article}

\usepackage[UTF8,scheme=plain,fontset=fandol]{ctex}
\usepackage[margin=22mm]{geometry}
\usepackage{amsmath,amssymb}
\usepackage[hidelinks]{hyperref}

\DeclareMathOperator{\softmax}{softmax}
\DeclareMathOperator{\diag}{diag}
\newcommand{\R}{\mathbb{R}}
\setlength{\parindent}{0pt}
\setlength{\parskip}{5pt}
\setlength{\emergencystretch}{2em}
\allowdisplaybreaks[1]

\begin{document}

\begin{center}
  {\Large\bfseries Assignment 2: AR and Transformers}\\[8pt]
  Name: 张逸林
  \hspace{8mm}
  Student ID: 2500013576
\end{center}

\section*{1.}

\paragraph{(a)}
Unmasked attention allows the model to attend to future tokens during
training and learn shortcut solutions, such as copying the target token
\(x_{i+1}\) instead of predicting it from \(x_1,\ldots,x_i\).
This violates the autoregressive constraint. Since future tokens are
unavailable during generation, it also creates a mismatch between training
and inference.

\paragraph{(b)}
Before applying the row-wise softmax, we set the attention scores at
positions \(j>i\) to \(-\infty\), while leaving the scores at positions
\(j\le i\) unchanged.
Since \(\exp(-\infty)=0\), the corresponding attention weights become zero
after softmax. Therefore, position \(i\) can attend only to the current and
previous tokens, preventing it from attending to future tokens.

\section*{2.}

\paragraph{(a)}
\[
Q=XW_Q,\qquad K=XW_K,\qquad V=XW_V,\qquad
\mathrm{output}=\softmax\!\left(\frac{QK^\top}{\sqrt d}\right)V.
\]
For the intermediate matrices, we have:
\[
\begin{aligned}
Q,K,V&\in\R^{n\times d},&
S&=\frac{QK^\top}{\sqrt d}\in\R^{n\times n},\\
A&=\softmax(S)\in\R^{n\times n},&
\mathrm{output}&=AV\in\R^{n\times d},
\end{aligned}
\]
where softmax is applied along each row.

\paragraph{(b)}
The time complexity of \(Q,K,V\) projections is
\[
3\cdot O(n\times d\times d)=O(nd^2).
\]
The time complexity of the attention computation is
\[
O(n^2d)+O(n^2)+O(n^2d)
=O\!\left(n^2(2d+1)\right)=O(n^2d).
\]
Thus, the overall time complexity is \(O(nd^2+n^2d)\).

The attention matrix has shape \(n\times n\) and contains \(n\times n\) elements.
Therefore, the memory complexity of explicitly storing the attention
matrix is \(O(n^2)\).

\clearpage
\section*{3.}
\begingroup
\small

Let
\[
q_i=(q_1,\ldots,q_{2h})^\top,\qquad
k_j=(k_1,\ldots,k_{2h})^\top.
\]
Let
\[
f(q_i,i)=[a_1,\ldots,a_{2h}]^\top,\qquad
f(k_j,j)=[b_1,\ldots,b_{2h}]^\top,\qquad
\theta=[\theta_1,\ldots,\theta_h]^\top,
\]
where \(\theta_x=b^{-2(x-1)/d}\), \(x=1,\ldots,h\).
By definition, we have:
\[
a_x=
\begin{cases}
q_x\cos(i\theta_x)-q_{h+x}\sin(i\theta_x),
& 1\le x\le h,\\
q_x\cos(i\theta_{x-h})+q_{x-h}\sin(i\theta_{x-h}),
& h+1\le x\le 2h,
\end{cases}
\]
\[
b_x=
\begin{cases}
k_x\cos(j\theta_x)-k_{h+x}\sin(j\theta_x),
& 1\le x\le h,\\
k_x\cos(j\theta_{x-h})+k_{x-h}\sin(j\theta_{x-h}),
& h+1\le x\le 2h.
\end{cases}
\]
Then we have:
\begin{align*}
f(q_i,i)^\top f(k_j,j)
&=\sum_{x=1}^{2h}a_xb_x\\
&=\sum_{x=1}^{h}
  \bigl(q_x\cos(i\theta_x)-q_{h+x}\sin(i\theta_x)\bigr)
  \bigl(k_x\cos(j\theta_x)-k_{h+x}\sin(j\theta_x)\bigr)\\
&\quad+\sum_{x=h+1}^{2h}
  \bigl(q_x\cos(i\theta_{x-h})+q_{x-h}\sin(i\theta_{x-h})\bigr)\\
&\hspace{37mm}\cdot
  \bigl(k_x\cos(j\theta_{x-h})+k_{x-h}\sin(j\theta_{x-h})\bigr)\\
&=\sum_{x=1}^{h}\Bigl[
 q_xk_x\cos(i\theta_x)\cos(j\theta_x)
 -q_xk_{h+x}\cos(i\theta_x)\sin(j\theta_x)\\
&\hspace{21mm}
 -q_{h+x}k_x\sin(i\theta_x)\cos(j\theta_x)
 +q_{h+x}k_{h+x}\sin(i\theta_x)\sin(j\theta_x)
 \Bigr]\\
&\quad+\sum_{x=h+1}^{2h}\Bigl[
 q_xk_x\cos(i\theta_{x-h})\cos(j\theta_{x-h})
 +q_xk_{x-h}\cos(i\theta_{x-h})\sin(j\theta_{x-h})\\
&\hspace{21mm}
 +q_{x-h}k_x\sin(i\theta_{x-h})\cos(j\theta_{x-h})
 +q_{x-h}k_{x-h}\sin(i\theta_{x-h})\sin(j\theta_{x-h})
 \Bigr]\\
&=\sum_{x=1}^{h}\Bigl[
 q_xk_x\bigl(\cos(i\theta_x)\cos(j\theta_x)
             +\sin(i\theta_x)\sin(j\theta_x)\bigr)\\
&\hspace{21mm}
 +q_xk_{h+x}\bigl(\sin(i\theta_x)\cos(j\theta_x)
                  -\sin(j\theta_x)\cos(i\theta_x)\bigr)\\
&\hspace{21mm}
 -q_{h+x}k_x\bigl(\sin(i\theta_x)\cos(j\theta_x)
                  -\sin(j\theta_x)\cos(i\theta_x)\bigr)\\
&\hspace{21mm}
 +q_{h+x}k_{h+x}\bigl(\cos(i\theta_x)\cos(j\theta_x)
                      +\sin(i\theta_x)\sin(j\theta_x)\bigr)
 \Bigr]\\
&=\sum_{x=1}^{h}\Bigl[
 q_xk_x\cos((i-j)\theta_x)
 +q_xk_{h+x}\sin((i-j)\theta_x)\\
&\hspace{21mm}
 -q_{h+x}k_x\sin((i-j)\theta_x)
 +q_{h+x}k_{h+x}\cos((i-j)\theta_x)
 \Bigr]\\
&=\sum_{x=1}^{2h}\sum_{y=1}^{2h}q_xk_y m_{xy}.
\end{align*}
Let \(\delta=i-j\). Define \(M(\delta)\in\R^{d\times d}\) as follows:
\[
M(\delta)=(m_{xy})_{x,y=1}^{2h}=
\begin{bmatrix}
\diag\!\bigl(\cos(\delta\theta)\bigr)
&
\diag\!\bigl(\sin(\delta\theta)\bigr)\\
-\diag\!\bigl(\sin(\delta\theta)\bigr)
&
\diag\!\bigl(\cos(\delta\theta)\bigr)
\end{bmatrix}.
\]
We obtain
\[
f(q_i,i)^\top f(k_j,j)=q_i^\top M(i-j)k_j.
\]
Thus, \(M(i-j)\) is independent of \(q_i\) and \(k_j\), and depends on the
positions only through \(i-j\).
\endgroup

\clearpage
\section*{4.}

The main idea of FlashAttention is to compute
exact attention faster and with less memory by reducing data transfers
between HBM and SRAM. It uses an IO-aware algorithm that splits \(Q\),
\(K\), and \(V\) into blocks and processes them in fast on-chip SRAM.
By incrementally updating the softmax statistics and the output, it
avoids storing the full attention matrix in HBM and repeatedly reading
and writing it, while preserving the result of standard attention.

FlashAttention is faster because it reduces memory transfers between
HBM and SRAM. It processes blocks of \(Q\), \(K\), and \(V\) in fast on-chip SRAM and fuses
attention operations into a single GPU kernel, avoiding repeated HBM
reads and writes of large intermediate matrices. Although recomputation
during the backward pass introduces additional arithmetic operations,
the reduction in memory traffic can more than compensate for this cost.

FlashAttention uses less memory because it does not store the full
\(n\times n\) attention matrix in HBM. Instead, it
incrementally updates the output and maintains per-row softmax
statistics. During the backward pass, the required attention blocks are
recomputed from the saved inputs and statistics. This reduces the attention
memory requirement from \(O(n^2+nd)\) to \(O(nd)\), which is linear in the sequence
length for fixed \(d\).

\end{document}
